\section{The Kolmogorov three-series theorem}
\label{sec-k-3-series}

As a further application of the Lindeberg-Feller
central limit theorem, we prove the Kolmogorov
three-series theorem (\hyperref[thm-three-series]{Theorem~\ref*{thm-three-series}}).

\begin{proof}[Proof of \texorpdfstring{\hyperref[thm-three-series]{Theorem~\ref*{thm-three-series}}}{Theorem}]
Denote $Y_n = X_n \mathbf{1}_{\{|X_n|\le 1\}}$.
Assume that the eponymous three series

\begin{align*}
\tag{17}\label{eq:k-first-series} &\sum_{n=1}^\infty \mathrm{P}(|X_n|>1) = \sum_{n=1}^\infty \mathrm{P}(X_n\neq Y_n), \\
\tag{18}\label{eq:k-second-series} &\sum_{n=1}^\infty \mathbf{E}(Y_n), \\
\tag{19}\label{eq:k-third-series} &\sum_{n=1}^\infty \mathbf{V}(Y_n)
\end{align*}

\noindent all converge. Applying \hyperref[thm-one-series]{Theorem~\ref*{thm-one-series}} to
the random variables $Y_n-\mathbf{E}(Y_n)$, from the
assumption that \eqref{eq:k-third-series} converges we
infer that the series
$\sum_{n=1}^\infty \big( Y_n-\mathbf{E}(Y_n) \big)$
converges almost surely. Since \eqref{eq:k-second-series}
converges, this also means that $\sum_n Y_n$
converges almost surely. Since \eqref{eq:k-first-series}
converges we get using the first Borel-Cantelli lemma
that the event $\{ X_n \neq Y_n\textrm{ i.o.} \}$ has
probability~$0$. Outside this event, the series
$\sum_n X_n$ converges if and only if $\sum_n Y_n$
converges, so $\sum_n X_n$ also converges almost
surely.

Next, for the converse claim, assume that $\sum_n X_n$
converges with positive probability. By the Kolmogorov
$0$-$1$ law, it therefore converges a.s. We need to
show that the series \eqref{eq:k-first-series},
\eqref{eq:k-second-series}, \eqref{eq:k-third-series}
converge.

First, note that \eqref{eq:k-first-series} must converge,
since if it doesn't, the second Borel-Cantelli lemma
implies that the event $\{ |X_n|>1 \textrm{ i.o.} \}$
has probability one,which would imply that $\sum X_n$
diverges almost surely.

Second, assume by contradiction that
\eqref{eq:k-third-series} diverges, or equivalently that

\[
v_n := \sum_{k=1}^n \mathbf{V}(Y_n) \to \infty\ \textrm{ as }n\to\infty.
\]

\noindent We apply \hyperref[thm-lindeberg-feller]{Theorem~\ref*{thm-lindeberg-feller}} to the
triangular array

\[
X_{n,k} = v_n^{-1/2} \left(Y_k - \mathbf{E} Y_k\right).
\]

\noindent Assumptions 1, 2 and 3 of the theorem are trivially
satisfied (with
$\sigma^2 = \sum_{k=1}^n \mathbf{V}(X_{n,k}) = 1$ for
all $n$). For assumption 4, note that
$|X_{n,k}| \le 2 v_n^{-1/2}$, so
$X_{n,k}^2 \mathbf{1}_{\{|X_{n,k}|>\epsilon\}} = 0$
for all $k$ if $v_n > 4/\epsilon^2$, which holds for
$n$ large enough. Thus, the assumptions of the
theorem are valid and hence we obtain the conclusion
that $S_n=\sum_{k=1}^n X_{n,k} \implies N(0,1)$ as
$n\to\infty$. Define a sequence (of numbers, not
random variables)

\[
t_n = v_n^{-1/2} \sum_{k=1}^n \mathbf{E} Y_k = S_n - v_n^{-1/2}\sum_{k=1}^n Y_k,
\]

\noindent and note that on the event
$\{\sum_n X_n \textrm{ converges} \}$ (in which case
clearly $\sum_n Y_n$ also converges), which we
assumed has probability~1,
$S_n-t_n = v_n^{-1/2} \sum_{k=1}^n Y_k \to 0$ as
$n\to\infty$. This is easily seen to be in
contradiction to the convergence in distribution of
$S_n$ to a limiting distribution of a non-constant
random variable. So we have shown that
\eqref{eq:k-third-series} converges.

Finally, having shown that \eqref{eq:k-third-series}
converges, we conclude from \hyperref[thm-one-series]{Theorem~\ref*{thm-one-series}}
that the series
$\sum_n \left(Y_n-\mathbf{E}(Y_n)\right)$ converges
almost surely. We already saw that on the event where
$\{ \sum_n X_n \textrm{ converges}\}$, which we
assumed has positive probability, also $\sum_n Y_n$
converges, and therefore the series
$\sum_n \mathbf{E}(Y_n) = \sum \big(Y_n -
(Y_n-\mathbf{E}(Y_n))\big)$
also converges.
\end{proof}
